% !TEX root = ../notes.tex

\documentclass[../notes.tex]{subfiles}

\begin{document}

\section{April 29}
Today, we start doing some number theory.

\subsection{Diamonds}
Let's recall our notation from last time. Let $k$ be $\FF_{p^{h(h-1)}}$ for some $h$. Recall that we are interested in $\mathbb B\coloneqq E(k,\Gamma_{h-1})\otimes_{W(k)}L_{K(h-1)}E(k,\Gamma_h)$, which is alternatively $E(L,\widetilde\Gamma_{n-1})$, where $\widetilde\Gamma_{n-1}$ is the universal deformation. Here, $L$ is an algebra over $K\coloneqq\FF_p((v_{n-1}))$, classifying isomorphisms $\widetilde\Gamma_{h-1}\cong\widetilde\Gamma_h^\circ$. It turns out that $L$ is not complete with respect to the $v_{n-1}$-adic topology, but we can complete it to $\widehat L$ and define $\widehat{\mathbb B}\coloneqq E(\widehat L,\widetilde\Gamma_h^\circ)$.

The action of the groups $\mathbb G_h$ and $\mathbb G_{h-1}$ on $\mathbb B$ produces $L_{K(h-1)}L_{K(h)}\mathbb S$ on taking fixed points. Our goal is to prove that $\widetilde{\mathbb B}$ has fixed points related to $L_{K(h-1)}\mathbb S$, so the ``wrong way map'' $\mathbb B\to\widehat{\mathbb B}$ should be able to prove a case of \Cref{conj:chromatic-splitting}. The Galois part of $\mathbb G_h$ turns out to be fairly unimportant, so we will focus on instead computing the fixed points for $\OO_h^\times$.
\begin{notation}
	For each $n$, we set $G_n\coloneqq\OO_n^\times$, where $\OO_n\coloneqq W(k)\langle S\rangle/\left(S^n-p,Sa-a^\sigma S\right)$.
\end{notation}
Accordingly, our goal for this chapter is to compute $\mathrm H^*_{\mathrm{cts}}(G_h^1;\widehat L)$.
\begin{remark}
	Last time, we asserted that $\widehat L\cong\FF_p((t^{1/p^\infty}))$, but this isomorphism is not helpful because it tells us nothing about the action by $G_h^1$!
\end{remark}
This will be done using the theory of diamonds. Diamonds are defined by gluing affinoid perfectoids together, so we should define those first.
\begin{definition}[power bounded]
	Let $R$ be a topological ring. Then we say that $x\in R$ is \textit{power bounded} if and only if the set of powers of $x$ is bounded. We let $R^\circ$ be the set of power bounded elements.
\end{definition}
\begin{definition}[affinoid perfectoid]
	Fix a positive characteristic $p$. An \textit{affinoid perfectoid} is a pair $(R,R^+)$ of classical topological rings satisfying the following.
	\begin{listalph}
		\item $R$ is a complete Hausdorff topological ring of characteristic $p$.
		\item Perfect: the Frobenius on $R$ is an isomorphism.
		\item The subring $R^+$ is open, integrally closed, and inside $R^\circ$.
		\item There exists a uniformizer $\varpi\in (R^+)^\times$ such that $R=R^+[1/\varpi]$, and $R^+$ has the $\varpi$-adic topology.
		\item The element $\varpi$ is topologically nilpotent unit.
	\end{listalph}
\end{definition}
\begin{remark}
	One can introduce condensed mathematics to avoid talking about topology, but we will not do so.
\end{remark}
\begin{example}
	In the sequel, we will take $R=\widehat L=\FF_p((t^{1/p^\infty}))$. Here, $R^+$ is $\FF_p[[t^{1/p^\infty}]]$, and our uniformizer is $t$.
\end{example}
\begin{remark}
	Technically, (d) implies (e) by the definition of the $\varpi$-adic topology.
\end{remark}
\begin{remark}
	The ring $R^+$ is perfect: indeed, for each $y\in R^+$, the polynomial $T^p-y$ admits a root in $R$, and this polynomial is integral over $R^+$, so this root must be in $R^+$.
\end{remark}
To glue, we need some maps.
\begin{definition}
	A morphism $f\colon(R,R^+)\to(S,S^+)$ of affinoid perfectoids is a continuous ring homomorphism $f\colon R\to S$ such that $f(R^+)\subseteq S^+$.
\end{definition}
\begin{remark}
	Equivalently, we could uniquely determine $f$ to be a continuous ring homomorphism $f^+\colon R^+\to S^+$ for which $\varpi$ goes to a topologically nilpotent unit in $S$.
\end{remark}
We are now ready to glue.
\begin{definition}[diamond]
	Fix a positive characteristic $p$. A \textit{diamond} is a functor from the category of affinoid perfectoids of characteristic $p$ to groupoids satisfying the $v$-sheaf condition, along with a nice cover.
\end{definition}
It will turn out that passing from a scheme to a diamond is enough to compute cohomology, which is the purpose of this theory for us.
\begin{remark}
	In the future, one hopes to generalize this perfectoid business to work for Gelfand rings, which already $L$ is an example of.
\end{remark}
\begin{example}
	There is a diamond $\op{Spd}\widehat L$, which is the functor represented by $\widehat L$.
\end{example}
Of course, we are interested in taking fixed points by $G_h^1$, so we want to make sense of $\op{Spd}(\widehat L)/G_h^1$. Then the cohomology of its structure sheaf will encode $\mathrm H^1(G_h^1;\widehat L)$, so we can compute it using perfectoid techniques.

\subsection{More Diamonds}
One can even find diamonds in characteristic zero, such as $\op{Spd}\QQ_p$. To make sense of the previous example, we need to explain tilting.
\begin{definition}[untilt]
	Fix an affinoid perfectoid $(R,R^+)$. An ideal $I$ in $W(R^+)$ is \textit{primitive of degree $1$} if and only if it is generated by a single element $\zeta=p+[\varpi]\alpha$, where $\alpha$ is a non-zero-divisor. For a choice of $I$ the \textit{untilt} of $(R,R^+)$ is the pair
	\[\left(W(R^+)[1/[\varpi]]/I,W(R^+)/I\right).\]
\end{definition}
\begin{remark}
	The tilt of this given untilt is $(R,R^+)$, where the tilt of $(S,S^+)$ is $S^+=\lim W(S^+)/(I,\varphi)$. It is a theorem that every untilt has one of the given forms.
\end{remark}
\begin{example}
	By definition, $\op{Spd}\QQ_p$ takes $(R,R^+)$ to the set of its untilts. Namely, a priori, one has the groupoid of pairs $(A,A^+)$ equipped with an isomorphism from the tilt of $(A,A^+)$ to $(R,R^+)$. It turns out that this groupoid is a set, which one checks by classifying untilts.
\end{example}
Here is an interesting example.
\begin{example}[Fargue--Fontaine curve]
	Fix a base affinoid perfectoid $(R,R^+)$; set $S\coloneqq\op{Spd}(R,R^+)$. Then we define the diamond $Y_S\coloneqq S\times\op{Spd}\QQ_p$, which sends $(A,A^+)$ to the collection of morphisms $(R,R^+)\to(A,A^+)$, along with an untilt of $(A,A^+)$. Now, for each power $n$ of $p$, define
	\[B_{S,[0,n]}\coloneqq W(R^+)\left\langle\frac{p^n}{[\varpi]}\right\rangle\left[\frac1{[\varpi]^{-1}}\right],\]
	and define $B_{S,[0,n]}^+$ to be the integral closure before taking the inverse. It turns out that
	\[Y_S=\left(\bigcup_n\op{Spa}\left(B_{S,[0,n]},B_{S,[0,n]}^+\right)\right)^\diamond.\]
	(We will not explain what $\op{Spa}$ is; it's just some adic space. We will also not explain what $(-)^\diamond$ is, which basically restricts a functor to affinoid perfectoids.) There is an action by Frobenius sending $B_{S,[0,n]}\to B_{S,[0,pn]}$, and we define $X_S$ to be the quotient of $Y_{S,(0,\infty)}$ by the Frobenius action.
\end{example}
We will eventually describe vector bundles over the Fargue--Fontaine curve.
\begin{definition}[Dieudonn\'e module]
	Fix a perfect ring $R$ of characteristic $p$. Then a \textit{Dieudonn\'e module} $M$ is a finite projective $W(R)$-module $M$ equipped with a Frobenius $F\colon M\otimes_{W(R),\varphi}W(R)\to M$ and Verschiebung $V\colon M\to M\otimes_{W(R),\varphi}W(R)$ satisfying $FV=VF=p$. Here, $-\otimes_{W(R),\varphi}W(R)$ means that the first $W(R)$ acts on the second through the Frobenius $\varphi$.
\end{definition}
\begin{remark}
	It turns out that $p$-divisible groups over $R$ are equivalent to Dieudonn\'e modules. Thus, we can study $p$-divisible groups by passing to linear algebra. Approximately speaking, this is akin to passing to the Lie algebra.
\end{remark}
\begin{remark}
	Given a Dieudonn\'e modules $M$, we see that $M[1/p]$ is a vector space over $W(R)[1/p]$, and the morphisms $F$ and $V$ extend to Frobenius semilinear automorphisms. Thus, if $(R,R^+)$ is affinoid perfectoid, a $p$-divisible group over $R^+$ gives rise to a quasicoherent sheaf on $X_{\op{Spd}(R,R^+)}$. The point is that the Dieudonn\'e module is already a module for the $B_{S,[0,n]}$s, and the fact that $F$ and $V$ are automorphisms allows us to descend to the quotient $X_S$ of $Y_S$ by the Frobenius.
\end{remark}
One may try to understand these Dieudonn\'e modules by understanding the induced quasicoherent sheaf, but this is not the full story because there is some integral structure that we forgot about.

Here is the relevance of the Fargue--Fontaine curve.
\begin{theorem}
	Fix an affinoid perfectoid $(R,R^+)$, and set $S\coloneqq\op{Spd}(R,R^+)$. Then $\op{Spd}\widehat L(S)$ classifies surjective maps $\OO(1/h)\to\OO(1/(h-1))$ of two vector bundles on $X_S$.
\end{theorem}
\begin{remark}
	Here, $\OO(1/h)$ is induced from a Dieudonn\'e module, which then is induced from a $p$-divisible group; namely, it comes from $\Gamma_h$. It has rank $h-1$.
\end{remark}
\begin{remark}
	By taking the kernel and twisting, it is equivalent to classify exact sequences
	\[0\to\OO\to\OO\left(\frac1h\right)\to\OO\left(\frac1{h-1}\right)\to0\]
	which preserve a choice of identification $\det\OO(1/h)\cong\det\OO(1/(h-1))$.
\end{remark}
We can now take our quotient. Viewing $G_h$ as the automorphisms of the Honda formal group, we receive a map $\det\colon G_h\to G_1$ given by taking the determinant of the corresponding Dieudonn\'e modules.
\begin{corollary}
	The quotient $\op{Spd}(\widehat L)/G_h^1$ classifies exact sequences
	\[0\to\OO\to\mc E\to\OO\left(\frac1{h-1}\right)\to0,\]
	where $\mc E$ is isomorphic to $\OO(1/h)$ (but without a chosen isomorphism).
\end{corollary}
\begin{remark}
	There is a diamond which classifies extensions on $X_S$ of $\OO$ by $\OO(1/(h-1))$, denoted $\op{BC}(-1/(h-1))$, where $\op{BC}$ stands for Banach--Colmez. We would like to study $\op{BC}(-1/(h-1))^*$, where $(-)^*$ means that we want the extension to be nontrivial. The Banach--Colmez spaces have well-understood cohomology, so our next task is to understand the cohomology in the punctured setting.
\end{remark}

\end{document}